\section{Preguntas probables de defensa y respuestas técnicas}

\subsection*{¿Por qué 114.000 filas no son 114.000 canciones?}
\textbf{Respuesta corta.} Porque solo hay 89.741 \code{track_id} únicos en el origen.\par
\textbf{Explicación técnica.} Una canción puede aparecer en varias filas por sus asignaciones de
género. Hay 16.641 IDs repetidos y 24.259 apariciones adicionales.\par
\textbf{Error que evitar.} Llamar ``duplicados exactos'' a esas filas: contienen distintos géneros y
no son idénticas.

\subsection*{¿Por qué consolidaron \code{track_id}?}
\textbf{Respuesta corta.} Para usar una canción como unidad y evitar doble conteo y fuga.\par
\textbf{Explicación técnica.} Si el mismo ID llega a entrenamiento y prueba, el desempeño queda
inflado. La consolidación mantiene atributos invariantes y reúne etiquetas.\par
\textbf{Error que evitar.} Eliminar todos los tracks repetidos y perder canciones válidas.

\subsection*{¿Por qué utilizaron la mediana de popularidad?}
\textbf{Respuesta corta.} Era una regla robusta, determinista y no dependía del orden de filas.\par
\textbf{Explicación técnica.} En 720 IDs repetidos el objetivo difería. Elegir el primero sería
arbitrario y la mediana reduce influencia de discrepancias extremas.\par
\textbf{Error que evitar.} Decir que la mediana recupera la popularidad ``verdadera'' o actual.

\subsection*{¿Qué es una correlación?}
\textbf{Respuesta corta.} Una medida de asociación entre variables.\par
\textbf{Explicación técnica.} Resume dirección e intensidad bajo una definición concreta; aquí se
usan asociaciones lineales y monotónicas. Sus valores van de $-1$ a 1.\par
\textbf{Error que evitar.} Interpretar un valor bajo como ausencia total de cualquier relación.

\subsection*{¿Qué diferencia existe entre Pearson y Spearman?}
\textbf{Respuesta corta.} Pearson mide asociación lineal; Spearman, asociación monotónica por
rangos.\par
\textbf{Explicación técnica.} Spearman tolera mejor curvaturas monotónicas y reduce sensibilidad a
magnitudes extremas, aunque ambos pueden fallar ante relaciones complejas.\par
\textbf{Error que evitar.} Afirmar que Spearman demuestra causalidad o resuelve todos los outliers.

\subsection*{¿Por qué correlación no implica causalidad?}
\textbf{Respuesta corta.} Porque una asociación puede surgir por confusores, selección o causalidad
inversa.\par
\textbf{Explicación técnica.} No hay intervención, temporalidad ni control suficiente de marketing,
género, recencia y exposición.\par
\textbf{Error que evitar.} Convertir un coeficiente en una recomendación para modificar la música.

\subsection*{¿Por qué \code{instrumentalness} no significa que una canción será impopular?}
\textbf{Respuesta corta.} Su correlación es débil y describe promedios de esta muestra.\par
\textbf{Explicación técnica.} Pearson es $-0{,}128$ y Spearman $-0{,}124$; existe amplio solapamiento
y posibles diferencias por género y exposición.\par
\textbf{Error que evitar.} Tratar una tendencia poblacional débil como destino de una canción.

\subsection*{¿Qué significa multicolinealidad?}
\textbf{Respuesta corta.} Que varios predictores contienen información lineal redundante.\par
\textbf{Explicación técnica.} \code{energy} correlaciona 0,759 con \code{loudness} y $-0{,}733$ con
\code{acousticness}. En modelos lineales puede desestabilizar coeficientes.\par
\textbf{Error que evitar.} Suponer que siempre obliga a eliminar una variable o que impide predecir.

\subsection*{¿Por qué no eliminaron outliers?}
\textbf{Respuesta corta.} Un extremo estadístico no es automáticamente un error.\par
\textbf{Explicación técnica.} Canciones largas, habladas, instrumentales o en vivo pueden ser casos
válidos. IQR se utilizó para revisar, y la fila inequívocamente inválida sí se retiró.\par
\textbf{Error que evitar.} Aplicar una regla IQR sin considerar el significado de la variable.

\subsection*{¿Qué problema tiene usar \code{artists} como predictor?}
\textbf{Respuesta corta.} Puede memorizar reputación y no generalizar a artistas nuevos.\par
\textbf{Explicación técnica.} Hay 31.437 valores únicos y un mismo artista puede aparecer en ambos
conjuntos. Se necesita control de encoding y validación agrupada.\par
\textbf{Error que evitar.} Presentar reconocimiento de identidades como aprendizaje musical general.

\subsection*{¿Qué problema tiene usar \code{track_id}?}
\textbf{Respuesta corta.} Es un identificador, no una característica generalizable.\par
\textbf{Explicación técnica.} Tiene un valor por canción procesada; un modelo puede memorizar y no
aprender un patrón reutilizable. Se conserva para trazabilidad y particiones.\par
\textbf{Error que evitar.} Codificarlo numéricamente o con one-hot como predictor.

\subsection*{¿Cómo funciona \code{track_genres}?}
\textbf{Respuesta corta.} Contiene etiquetas únicas separadas por \code{|}.\par
\textbf{Explicación técnica.} Al consolidar un track se preservan todos sus géneros ordenados. Para
analizar se expande; para modelar se crean indicadores por etiqueta.\par
\textbf{Error que evitar.} Contar las filas expandidas como canciones nuevas.

\subsection*{¿Qué es multi-hot encoding?}
\textbf{Respuesta corta.} Una columna binaria por etiqueta, permitiendo varias activas por fila.\par
\textbf{Explicación técnica.} Una canción \code{pop|rock} tendrá 1 en ambas columnas. A diferencia de
one-hot exclusivo, respeta pertenencia múltiple.\par
\textbf{Error que evitar.} Elegir solo el primer género y descartar información.

\subsection*{¿Qué sesgos tiene el dataset?}
\textbf{Respuesta corta.} Muestreo balanceado por género, temporalidad desconocida, variables
omitidas y dependencia por artista.\par
\textbf{Explicación técnica.} Estos sesgos afectan agregados, generalización temporal y evaluación
de artistas nuevos. La consolidación corrige duplicación, no representatividad.\par
\textbf{Error que evitar.} Afirmar que 1.000 filas por género reflejan el catálogo real.

\subsection*{¿Qué riesgos éticos existen?}
\textbf{Respuesta corta.} Reforzar popularidad previa, reducir diversidad y automatizar promoción.\par
\textbf{Explicación técnica.} Un sistema optimizado con señales históricas puede reproducir
desigualdades de exposición y convertir correlaciones en criterios normativos.\par
\textbf{Error que evitar.} Equiparar popularidad estimada con calidad o mérito artístico.

\subsection*{¿Qué privacidad hay que considerar?}
\textbf{Respuesta corta.} El riesgo actual es bajo porque no hay datos de oyentes, pero no es cero.\par
\textbf{Explicación técnica.} Existen metadatos públicos de obras y artistas; se exige trazabilidad y
uso responsable. Integrar perfiles o escuchas requeriría controles adicionales.\par
\textbf{Error que evitar.} Concluir que todo dato musical puede reutilizarse sin condiciones.

\subsection*{¿Por qué \code{popularity} es problemática?}
\textbf{Respuesta corta.} Es dinámica, histórica y depende de factores no observados.\par
\textbf{Explicación técnica.} Reproducciones y recencia influyen, no se conoce la fecha de extracción
y Spotify marca el campo como obsoleto en la API actual.\par
\textbf{Error que evitar.} Tratarla como una medida atemporal de éxito o calidad.

\subsection*{¿Qué harían distinto con datos temporales?}
\textbf{Respuesta corta.} Definirían un horizonte y validarían entrenando en pasado y probando en
futuro.\par
\textbf{Explicación técnica.} Se crearían variables disponibles al momento de predicción, evitando
información posterior, y se mediría deriva entre periodos.\par
\textbf{Error que evitar.} Mezclar fechas aleatoriamente y afirmar capacidad prospectiva.

\subsection*{¿Por qué usar MAE, RMSE y $R^2$?}
\textbf{Respuesta corta.} Miden aspectos complementarios del error.\par
\textbf{Explicación técnica.} MAE es interpretable en puntos; RMSE enfatiza errores grandes; $R^2$
compara variabilidad explicada contra una referencia. Deben calcularse fuera de muestra.\par
\textbf{Error que evitar.} Elegir una sola métrica o interpretar $R^2$ como porcentaje de aciertos.

\subsection*{¿La diferencia de \code{explicit} es importante?}
\textbf{Respuesta corta.} Es visible, pero pequeña en escala estandarizada.\par
\textbf{Explicación técnica.} La media difiere 4,03 puntos y $d$ de Cohen es 0,196. Con muestras
grandes puede existir significancia estadística sin gran relevancia práctica.\par
\textbf{Error que evitar.} Recomendar contenido explícito basándose en una comparación no causal.

\subsection*{¿Por qué las correlaciones débiles no descartan un modelo?}
\textbf{Respuesta corta.} Un modelo puede combinar señales, interacciones y curvaturas.\par
\textbf{Explicación técnica.} Las correlaciones son pares y resumen formas simples; el valor real debe
medirse fuera de muestra contra una línea base.\par
\textbf{Error que evitar.} Prometer que complejidad garantiza buen desempeño.
